\section{Extension: Signature Schemes and Long-Lived Evidence}
\label{sec:ext-schemes}
% Self-contained extension: no invariant of the core is changed or weakened.
% Written to move unchanged into the extension part of a core/extension split.
% ================================================================================

This section is an extension. It adds conventions and guidance and removes
no MUST: every invariant of \S\ref{sec:additional-invariants} holds with
any scheme whose verifier is unforgeable, because the contract constrains
\emph{what a signature binds} (the domain-separated tuple of
Appendix~\ref{app:wire-format}), not how it is computed.

\paragraph{Scheme tag and registry convention.} Each client key is
registered with a scheme tag and verification material (I15), and the
tag is part of the signed registration tuple, so a key cannot be
re-read under another scheme. A deployment keeps a scheme list: a mapping
from tag to a verifier and to the encoding of its public key and
signature. The reference list holds \texttt{ed25519} and
\texttt{ecdsa-p256}. A tag not on the list is invalid, never a fallback to
another scheme, and a verifier answers false for malformed material. A
further scheme, a post-quantum one such as ML-DSA, or a hybrid that
requires both signatures to verify, is added by registering a tag and a
verifier; the bridge, the state machine and the log format are unchanged.
Tags SHOULD name the algorithm and parameter set, and a hybrid SHOULD have
its own tag rather than reuse a component's. The schemas list the
reference tags and are extended by the deployment with its own. Evidence:
a test-only stub scheme is accepted through the same verifier on both
runners (T144, and the test \texttt{a\_further\_scheme\_is\_accepted\_through\_the\_verifier}
of the Rust crate), and the same key is refused where the verifier lacks
the tag; it exercises dispatch only and says nothing about a real
post-quantum scheme.

\paragraph{Long-lived signatures.} The signatures differ in how long they
must hold. A re-attach, join or key-addition signature matters while its
nonce is live and, for the key, while the session can still be resumed,
so migrating client keys to a new scheme is a rotation (I15,
I11-REVOKE) that needs no change to past events. The operator
signatures on events and checkpoints are archived evidence and outlive any
algorithm. For these:
\begin{itemize}
\item Sign checkpoints (I-CHECKPOINT) with a post-quantum or hybrid scheme
 first; a checkpoint covers every earlier event through the Merkle root,
 so protecting it protects the prefix.
\item Before a scheme is retired, re-sign: write a new checkpoint over the
 current head under the new scheme, and keep the old checkpoint. The
 hash chain, not the old signature, then carries integrity of the prefix, as
 long as the hash function remains collision resistant.
\item Anchor the checkpoint root outside the bridge (a transparency
 service or a timestamp) before the old algorithm weakens, so that
 evidence made earlier cannot be backdated after a break.
\item Record the scheme tag with every signature and allow several
 signatures over one checkpoint, so a verifier accepts evidence that
 verifies under any scheme it still trusts.
\end{itemize}

\paragraph{Algorithm agility without downgrade.} A verifier trusts a
configured list, not the tag a signer presents. Removing a scheme from the
list rejects new signatures under it, and archived evidence under it is
accepted only with an anchor or a later re-signature as above. The bridge
never negotiates a scheme with the client. Hash agility (the chain and
Merkle hash) is out of scope here.
